\section{Troubleshooting}

\subsection{Neo4j non si Avvia}

\textbf{Sintomo}: Il container di Neo4j mostra status \texttt{Exited}

\textbf{Soluzione}:
\begin{enumerate}
    \item Verifica le credenziali nel file \texttt{.env}
    \item Controlla lo spazio su disco disponibile
    \item Visualizza i log: \texttt{docker-compose logs neo4j}
    \item Riavvia: \texttt{docker-compose down -v \&\& docker-compose up -d}
\end{enumerate}

\subsection{Ollama Timeout}

\textbf{Sintomo}: Errori di timeout quando si chiama l'API Ollama

\textbf{Soluzione}:
\begin{enumerate}
    \item Scarica i modelli manualmente: \texttt{ollama pull qwen2.5:7b}
    \item Verifica la RAM disponibile (minimo 8 GB)
    \item Aumenta il timeout nelle configurazioni se necessario
\end{enumerate}

\subsection{Porta Già in Uso}

\textbf{Sintomo}: Errore ``Address already in use''

\textbf{Soluzione}: Modifica le porte in \texttt{docker-compose.yml}:
\begin{lstlisting}
ports:
  - "3001:3000"  # Frontend su porta 3001 invece di 3000
\end{lstlisting}

\subsection{Container Crasha Subito Dopo l'Avvio}

\textbf{Sintomo}: I container si fermano immediatamente

\textbf{Soluzione}:
\begin{enumerate}
    \item Visualizza i log: \texttt{docker-compose logs}
    \item Verifica la configurazione nel file \texttt{.env}
    \item Assicurati che il file \texttt{.env} sia nella directory radice
\end{enumerate}

\subsection{Nessuna Connessione al Database}

\textbf{Sintomo}: L'applicazione non si connette a Neo4j

\textbf{Soluzione}:
\begin{enumerate}
    \item Verifica che Neo4j sia in stato \texttt{healthy}: \texttt{docker-compose ps}
    \item Accedi a Neo4j Browser: \href{http://localhost:7474}{http://localhost:7474}
    \item Controlla la variabile \texttt{NEO4J\_URI} nel file \texttt{.env}
\end{enumerate}
